\chapter{Lezione 8 --- Anti-trasformata z, implementazione dei filtri digitali, trasformata bilineare di Tustin e pre-warping}

\section{Richiamo: tavola essenziale per l'anti-trasformata z}

Negli appunti viene ripresa una tabella di coppie utili tra segnale continuo causale,
successione campionata e trasformata z.

Ponendo
\[
\varphi=e^{-aT},
\]
si hanno le seguenti corrispondenze fondamentali:

\[
\mathcal{Z}\{1\}=\frac{z}{z-1},
\]
\[
\mathcal{Z}\{kT\}=\frac{Tz}{(z-1)^2},
\]
\[
\mathcal{Z}\left\{\frac{1}{2}k^2T^2\right\}
=
\frac{T^2}{2}\frac{z(z+1)}{(z-1)^3},
\]
\[
\mathcal{Z}\{\varphi^k\}=\frac{z}{z-\varphi},
\]
\[
\mathcal{Z}\{kT\,\varphi^k\}
=
\frac{T\varphi z}{(z-\varphi)^2},
\]
\[
\mathcal{Z}\{1-\varphi^k\}
=
\frac{z(1-\varphi)}{(z-1)(z-\varphi)},
\]
\[
\mathcal{Z}\{\sin(\omega kT)\}
=
\frac{z\sin(\omega T)}{z^2-2\cos(\omega T)\,z+1},
\]
\[
\mathcal{Z}\{\cos(\omega kT)\}
=
\frac{z\bigl(z-\cos(\omega T)\bigr)}{z^2-2\cos(\omega T)\,z+1},
\]
\[
\mathcal{Z}\{\varphi^k\sin(\omega kT)\}
=
\frac{z\varphi\sin(\omega T)}{z^2-2\varphi\cos(\omega T)\,z+\varphi^2},
\]
\[
\mathcal{Z}\{\varphi^k\cos(\omega kT)\}
=
\frac{z\bigl(z-\varphi\cos(\omega T)\bigr)}{z^2-2\varphi\cos(\omega T)\,z+\varphi^2}.
\]

Questa tabella \`e il riferimento operativo per tornare da una funzione razionale \(G(z)\)
alla corrispondente successione nel tempo discreto.

\section{Anti-trasformata z con il metodo dei fratti semplici}

Negli appunti si osserva che, diversamente da quanto accade in Laplace,
nelle forme tabulate della trasformata z compare quasi sempre un fattore \(z\) al numeratore.
Per questo, data una funzione razionale
\[
G(z)=\frac{N(z)}{D(z)},
\]
per applicare il metodo dei fratti semplici conviene spesso considerare
\[
\frac{G(z)}{z}.
\]

L'idea operativa \`e:
\[
G(z)
\quad\longrightarrow\quad
\frac{G(z)}{z}
\quad\longrightarrow\quad
\text{espansione in fratti semplici}
\quad\longrightarrow\quad
\text{calcolo dei residui}
\quad\longrightarrow\quad
\text{moltiplico per } z
\quad\longrightarrow\quad
\text{anti-trasformo via tabella}.
\]

In questo modo si ottiene una soluzione analitica della successione.

\subsection*{Osservazione}

Il procedimento \`e del tutto analogo a quello usato nel dominio di Laplace:
si fattorizza il denominatore, si espande in termini elementari e poi si riconosce
l'anti-trasformata di ciascun termine.

\section{Metodo non analitico: lunga divisione}

Negli appunti viene richiamato anche un secondo metodo, pi\`u algoritmico:
la \textbf{lunga divisione} tra numeratore e denominatore.

Se
\[
G(z)=\frac{N(z)}{D(z)},
\]
si sviluppa la funzione come serie in potenze di \(z^{-1}\):
\[
G(z)=g(0)+g(1)z^{-1}+g(2)z^{-2}+\dots
\]

I coefficienti letti nella serie sono proprio i valori della successione.

\subsection*{Interpretazione}

Ogni moltiplicazione per \(z^{-1}\) corrisponde a un ritardo di un campione,
mentre ogni moltiplicazione per \(z\) corrisponde formalmente a un anticipo.

Il metodo della lunga divisione \`e particolarmente utile quando:
\begin{itemize}
    \item si vogliono i primi campioni della risposta;
    \item la fattorizzazione del denominatore \`e scomoda;
    \item interessa un approccio numerico pi\`u che una formula chiusa.
\end{itemize}

\section{Dalla funzione di trasferimento all'equazione alle differenze}

Consideriamo una funzione di trasferimento discreta del tipo
\[
G(z)=\frac{n(z)}{d(z)}
=
\frac{b_m z^m+\dots+b_0}{z^n+a_{n-1}z^{n-1}+\dots+a_0}.
\]

Assumendo la relazione ingresso--uscita
\[
Y(z)=G(z)U(z),
\]
moltiplicando per il denominatore si ottiene
\[
d(z)Y(z)=n(z)U(z).
\]

Anti-trasformando, si ricava una equazione alle differenze del tipo
\[
y(k)+a_{n-1}y(k-1)+\dots+a_0 y(k-n)
=
b_m u(k)+b_{m-1}u(k-1)+\dots+b_0 u(k-m).
\]

Quindi si pu\`o risolvere esplicitamente rispetto a \(y(k)\):
\[
y(k)=
-a_{n-1}y(k-1)-\dots-a_0y(k-n)
+b_m u(k)+\dots+b_0u(k-m).
\]

Questa \`e l'equazione che permette di implementare il filtro nel tempo discreto.

\section{Implementazione digitale del filtro}

Per implementare l'equazione alle differenze su una macchina digitale serve:
\begin{itemize}
    \item un \textbf{buffer di memoria}, per memorizzare i campioni passati di ingresso e uscita;
    \item una \textbf{ALU} (\emph{Arithmetic Logic Unit}), che esegue somme e moltiplicazioni;
    \item una logica di aggiornamento dei campioni nel buffer.
\end{itemize}

Se l'ordine del filtro \`e \(n\), in una realizzazione tipica bisogna conservare:
\begin{itemize}
    \item \(n\) campioni passati dell'uscita;
    \item \(n+1\) campioni dell'ingresso.
\end{itemize}

Quindi il buffer totale contiene, a meno di convenzioni di indicizzazione,
circa \(2n+1\) campioni.

\section{Filtri IIR e filtri FIR}

\subsection{Filtri IIR}

Un filtro descritto da una equazione alle differenze con dipendenza
anche dai valori passati dell'uscita \`e un filtro \textbf{ricorsivo}.

Questi filtri si chiamano
\[
\text{IIR} \quad \text{(\emph{Infinite Impulse Response})}.
\]

Il nome deriva dal fatto che, in generale, la risposta all'impulso discreto
non si annulla in un numero finito di passi.

\subsection{Filtri FIR}

Se invece la funzione di trasferimento ha tutti i poli nell'origine, cio\`e pu\`o essere scritta come
\[
G(z)=\frac{m(z)}{z^n},
\]
allora il filtro non \`e ricorsivo e l'uscita dipende solo da valori presenti e passati dell'ingresso:
\[
y(k)=b_m u(k)+b_{m-1}u(k-1)+\dots+b_0 u(k-m).
\]

Questi filtri si chiamano
\[
\text{FIR} \quad \text{(\emph{Finite Impulse Response})}.
\]

Se l'ingresso \(u(k)\) \`e un impulso discreto, la risposta di un filtro FIR va a zero
in un numero finito di passi.

\section{Specifiche di un filtro passa-basso}

Negli appunti si introduce il problema progettuale:
si vogliono assegnare le specifiche di un filtro passa-basso nel continuo
e poi costruire un filtro digitale equivalente.

Un modello analogico elementare \`e
\[
G(s)=\frac{\omega_b}{s+\omega_b},
\]
dove \(\omega_b\) \`e la pulsazione di banda.

Si tratta di un filtro del primo ordine, con pendenza asintotica
\[
-20\ \text{dB/decade}
\]
oltre la frequenza di taglio.

\section{Problema: come passare da \(s\) a \(z\)?}

Se si vuole lavorare non pi\`u su segnali continui ma su segnali campionati,
bisogna sostituire il filtro analogico con un filtro discreto \(G(z)\).

La relazione esatta tra \(s\) e \(z\),
\[
z=e^{sT},
\]
\`e per\`o trascendente, quindi non porta in generale a una funzione razionale semplice in \(z\).

Per ottenere una funzione di trasferimento digitale razionale
si introducono trasformazioni approssimate da \(s\) a \(z\).

\section{Idea generale: approssimazioni razionali}

Negli appunti si richiama che la trasformazione non pu\`o essere esatta,
ma si pu\`o cercare una buona approssimazione.

Una prima idea \`e quella di usare lo sviluppo di Taylor troncato.
Una scelta pi\`u efficace \`e usare un approssimante di Pad\'e,
cio\`e una funzione razionale che approssima una funzione trascendente.

Qui si sceglie un approssimante di Pad\'e del primo ordine per \(e^{sT}\).

\section{Approssimazione di Pad\'e del primo ordine}

Scriviamo
\[
e^{sT}= \frac{e^{\frac{sT}{2}}}{e^{-\frac{sT}{2}}}.
\]

Approssimando numeratore e denominatore al primo ordine di Taylor:
\[
e^{\frac{sT}{2}} \approx 1+\frac{sT}{2},
\qquad
e^{-\frac{sT}{2}} \approx 1-\frac{sT}{2}.
\]

Si ottiene quindi
\[
e^{sT}\approx \frac{1+\frac{sT}{2}}{1-\frac{sT}{2}}.
\]

Poich\'e \(z=e^{sT}\), segue l'approssimazione
\[
z\approx \frac{1+\frac{sT}{2}}{1-\frac{sT}{2}}.
\]

Risolvendo rispetto a \(s\):
\[
\left(1-\frac{sT}{2}\right)z = 1+\frac{sT}{2},
\]
\[
z-1=\frac{sT}{2}(z+1),
\]
\[
s=\frac{2}{T}\frac{z-1}{z+1}.
\]

Questa \`e la \textbf{trasformata di Tustin}, detta anche \textbf{trasformata bilineare}.

\section{Trasformata di Tustin o bilineare}

La sostituzione fondamentale \`e
\[
s=\frac{2}{T}\frac{z-1}{z+1}.
\]

Essa permette di trasformare una funzione di trasferimento continua \(G(s)\)
in una funzione di trasferimento discreta \(G(z)\) razionale.

\subsection*{Osservazione}

La trasformazione \`e semplice e molto usata in pratica,
ma resta comunque un'approssimazione del legame esatto tra \(s\) e \(z\).

\section{Esempio: discretizzazione di un passa-basso del primo ordine}

Consideriamo il filtro analogico
\[
G(s)=\frac{\omega_b}{s+\omega_b}.
\]

Sostituendo \(s=\frac{2}{T}\frac{z-1}{z+1}\), si ottiene
\[
G(z)=\frac{\omega_b}{\frac{2}{T}\frac{z-1}{z+1}+\omega_b}.
\]

Moltiplicando numeratore e denominatore per \(z+1\):
\[
G(z)=
\frac{\omega_b(z+1)}
{\frac{2}{T}(z-1)+\omega_b(z+1)}.
\]

Moltiplicando ancora per \(T\):
\[
G(z)=
\frac{T\omega_b(z+1)}
{2(z-1)+T\omega_b(z+1)}.
\]

Sviluppando il denominatore:
\[
G(z)=
\frac{T\omega_b(z+1)}
{(2+T\omega_b)z-2+T\omega_b}.
\]

Questa \`e una funzione razionale in \(z\), quindi pu\`o essere implementata
come filtro digitale tramite le tecniche gi\`a viste.

\section{Dove vanno i punti del piano \(s\) con la trasformata bilineare?}

La lezione si chiede anche quale sia l'effetto geometrico della trasformazione di Tustin
sui punti del piano \(s\).

Prendendo un punto sull'asse immaginario,
\[
s=j\omega,
\]
si ha
\[
z=\frac{1+j\omega T/2}{1-j\omega T/2}.
\]

Calcoliamo il modulo:
\[
|z|^2
=
\frac{\left(1+j\omega T/2\right)\left(1-j\omega T/2\right)}
{\left(1-j\omega T/2\right)\left(1+j\omega T/2\right)}
=1.
\]

Quindi
\[
|z|=1.
\]

Ne segue che l'asse immaginario del piano \(s\) viene mappato sul cerchio unitario del piano \(z\).

Inoltre:
\begin{itemize}
    \item il semipiano sinistro di \(s\) viene mappato all'interno del cerchio unitario;
    \item il semipiano destro di \(s\) viene mappato all'esterno del cerchio unitario.
\end{itemize}

Questa propriet\`a \`e molto importante perch\'e preserva il criterio di stabilit\`a:
un sistema stabile nel continuo viene trasformato in un sistema stabile nel discreto.

\section{Distorsione in frequenza: frequency warping}

La trasformata di Tustin non preserva linearmente la frequenza.
Per capire cosa succede alla risposta armonica,
si valuta il filtro discreto sul cerchio unitario:
\[
z=e^{j\omega_d T}.
\]

Allora
\[
G_d(e^{j\omega_d T})
=
G_a\!\left(\frac{2}{T}\frac{e^{j\omega_d T}-1}{e^{j\omega_d T}+1}\right).
\]

Semplificando il rapporto:
\[
\frac{e^{j\omega_d T}-1}{e^{j\omega_d T}+1}
=
\frac{e^{j\omega_d T/2}\left(e^{j\omega_d T/2}-e^{-j\omega_d T/2}\right)}
{e^{j\omega_d T/2}\left(e^{j\omega_d T/2}+e^{-j\omega_d T/2}\right)}
=
\frac{2j\sin(\omega_d T/2)}{2\cos(\omega_d T/2)}
=
j\tan\!\left(\frac{\omega_d T}{2}\right).
\]

Quindi
\[
G_d(e^{j\omega_d T})
=
G_a\!\left(
j\,\frac{2}{T}\tan\!\left(\frac{\omega_d T}{2}\right)
\right).
\]

Questo mostra che le pulsazioni analogiche e discrete sono legate da
\[
\omega_a=\frac{2}{T}\tan\!\left(\frac{\omega_d T}{2}\right),
\]
oppure inversamente
\[
\omega_d=\frac{2}{T}\arctan\!\left(\frac{\omega_a T}{2}\right).
\]

\subsection*{Interpretazione}

La risposta armonica viene distorta poco alle basse frequenze
e molto di pi\`u alle alte frequenze.

Per questo si parla di \textbf{warping} della frequenza:
la scala delle pulsazioni nel continuo non viene trasportata linearmente
nella scala discreta.

\section{Pre-warping}

Spesso interessa che il filtro discreto riproduca in modo esatto
il filtro analogico almeno a una specifica pulsazione \(\omega_0\),
ad esempio proprio alla frequenza di taglio.

Per ottenere questo, si modifica la trasformazione bilineare introducendo una costante \(k\):
\[
s=k\,\frac{z-1}{z+1}.
\]

Si sceglie \(k\) imponendo che la corrispondenza sia esatta alla pulsazione desiderata \(\omega_0\):
\[
\omega_0 = k\,\tan\!\left(\frac{\omega_0 T}{2}\right),
\]
da cui
\[
k=\frac{\omega_0}{\tan\!\left(\frac{\omega_0 T}{2}\right)}.
\]

Questa scelta prende il nome di \textbf{pre-warping}.

\subsection*{Significato}

Con il pre-warping si forza la trasformazione in modo che
la risposta armonica del filtro discreto coincida con quella del filtro analogico
alla frequenza \(\omega_0\) di progetto.

In altre parole:
\begin{itemize}
    \item senza pre-warping, la Tustin \`e semplice ma distorce la scala in frequenza;
    \item con pre-warping, si corregge localmente la distorsione nel punto di maggiore interesse.
\end{itemize}

\section{Sintesi della lezione}

I punti chiave della lezione sono:
\begin{enumerate}
    \item richiamo dei metodi per anti-trasformare una funzione razionale in \(z\):
    \begin{itemize}
        \item fratti semplici;
        \item lunga divisione;
    \end{itemize}
    \item da \(Y(z)=G(z)U(z)\) si ricava l'equazione alle differenze che implementa il filtro;
    \item i filtri ricorsivi sono filtri IIR, mentre quelli non ricorsivi con poli nell'origine
    sono filtri FIR;
    \item per discretizzare un filtro analogico si cerca una trasformazione razionale da \(s\) a \(z\);
    \item usando l'approssimante di Pad\'e del primo ordine si ottiene la trasformata di Tustin:
    \[
    s=\frac{2}{T}\frac{z-1}{z+1};
    \]
    \item applicando Tustin a
    \[
    G(s)=\frac{\omega_b}{s+\omega_b}
    \]
    si ottiene un filtro discreto razionale \(G(z)\);
    \item la trasformata bilineare manda l'asse immaginario di \(s\) sul cerchio unitario di \(z\);
    \item la corrispondenza tra frequenze \`e
    \[
    \omega_a=\frac{2}{T}\tan\!\left(\frac{\omega_d T}{2}\right),
    \]
    quindi la scala in frequenza \`e deformata;
    \item con il pre-warping si usa
    \[
    s=k\frac{z-1}{z+1},
    \qquad
    k=\frac{\omega_0}{\tan(\omega_0 T/2)},
    \]
    per garantire corrispondenza esatta alla frequenza \(\omega_0\).
\end{enumerate}